\section{E1: Erkennung der Verbrauchsänderungen}
\label{app:report_e1}
\subsection*{Setup}
\begin{itemize}
\item \textbf{Lauf:} e1\_\allowbreak{}006
\item \textbf{Stichtage:} SF6: 02.12.2016, SF7: 13.01.2017, SF8: 20.02.2017
\item \textbf{Szenarien:} scaled, scaled\_\allowbreak{}down, ramp\_\allowbreak{}28, ramp\_\allowbreak{}28\_\allowbreak{}down, ramp\_\allowbreak{}56, ramp\_\allowbreak{}56\_\allowbreak{}down, peak
\item \textbf{Systeme:} SYS1, die Regel ordnet ein, keine Adaption, kein Sprachmodell
\item \textbf{Commit:} f647560
\end{itemize}

\subsection*{Kriterien der Teilfrage 1}
\begin{tablang}{W{2.05} W{0.43} W{0.92} W{0.59}}{Kriterium & Ziel & Ergebnis & erfüllt}
Erkennungsrate eingebauter Anomalien & $\geq$ 90 \% & 9 von 13 (69 \%) & nein \\
Erkennungsrate langsamer Rampen & 0 \% & 3 von 6 (50 \%) & nein \\
\end{tablang}

\subsection*{Erkennung je Szenario über die Stichtage}
\begin{tablang}{W{0.88} W{0.94} W{0.69} W{1.25} W{1.25}}{Szenario & erwartet & erkannt & Erkennungstag, Mittel (Spanne) & auslösende Skala}
scaled & erkannt & 1 von 3 & 15,0 (15) & weekly + long\_\allowbreak{}term \\
scaled\_\allowbreak{}down & erkannt & 2 von 3 & 34,5 (31–38) & weekly + long\_\allowbreak{}term \\
ramp\_\allowbreak{}28 & erkannt & 2 von 3 & 33,5 (17–50) & weekly + long\_\allowbreak{}term \\
ramp\_\allowbreak{}28\_\allowbreak{}down & erkannt & 3 von 3 & 43,0 (31–57) & weekly + long\_\allowbreak{}term \\
ramp\_\allowbreak{}56 & nicht erkannt & 2 von 3 & 43,0 (29–57) & weekly + long\_\allowbreak{}term \\
ramp\_\allowbreak{}56\_\allowbreak{}down & nicht erkannt & 1 von 3 & 57,0 (57) & weekly + long\_\allowbreak{}term \\
peak & erkannt & 1 von 1 & - & weekday \\
\end{tablang}

\subsection*{Hinweise}
\begin{itemize}
\item Tag = Entscheidungstag seit dem Stichtag, an Tag 1 wird der Stichtag selbst eingeordnet
\item erkannt: Lastspitze = Einzelanomalie (Skala daily oder weekday) an Tag 1, Dauerlast und Rampe = anhaltende Veränderung (weekly und long\_\allowbreak{}term an 6 von 7 Tagen). Gezählt wird der erste Tag, an dem der Szenariolauf den Kennwert hat und der Kontrolllauf desselben Stichtags nicht
\item auswertbar: zeigt der Kontrolllauf am Einbautag der Lastspitze schon eine Einzelanomalie, geht der Fall nicht in die Zählung ein
\item erwartet: die Rampe über 56 Tage soll nicht als anhaltende Veränderung erkannt werden
\item Anhaltende Veränderungen werden im Mittel an Tag 35,0 erkannt (Dauerlast und Rampe über 28 Tage)
\item Erkennungstag, Mittel (Spanne): nur über die erkannten Fälle, bei der Lastspitze entfällt er, sie zählt nur an Tag 1
\item Skalen anhaltend: weekly bzw. long\_\allowbreak{}term liegt an 6 der letzten 7 Tage über der Schwelle, eine anhaltende Veränderung gilt erst, wenn beide gleichzeitig anhalten. Abgerissen zählt, wie oft weekly vor diesem Tag wieder unter 6 von 7 Tagen fiel
\end{itemize}

\subsection*{Weitere Tabellen}

\subsubsection*{Erkennung je Stichtag und Szenario}
\begin{tablang}{W{0.80} W{0.93} V{0.73} W{1.33} V{1.20} W{1.00}}{Stichtag & Szenario & erkannt am Tag & auslösende Skala & Tage mit Kennwert im Kontrolllauf & auswertbar}
SF6 & scaled & - & - & 0 & ja \\
SF6 & scaled\_\allowbreak{}down & 31 & weekly + long\_\allowbreak{}term & 0 & ja \\
SF6 & ramp\_\allowbreak{}28 & 50 & weekly + long\_\allowbreak{}term & 0 & ja \\
SF6 & ramp\_\allowbreak{}28\_\allowbreak{}down & 31 & weekly + long\_\allowbreak{}term & 0 & ja \\
SF6 & ramp\_\allowbreak{}56 & 57 & weekly + long\_\allowbreak{}term & 0 & ja \\
SF6 & ramp\_\allowbreak{}56\_\allowbreak{}down & - & - & 0 & ja \\
SF6 & peak & - & - & 1 & nein \\
SF7 & scaled & 15 & weekly + long\_\allowbreak{}term & 0 & ja \\
SF7 & scaled\_\allowbreak{}down & 38 & weekly + long\_\allowbreak{}term & 0 & ja \\
SF7 & ramp\_\allowbreak{}28 & 17 & weekly + long\_\allowbreak{}term & 0 & ja \\
SF7 & ramp\_\allowbreak{}28\_\allowbreak{}down & 41 & weekly + long\_\allowbreak{}term & 0 & ja \\
SF7 & ramp\_\allowbreak{}56 & 29 & weekly + long\_\allowbreak{}term & 0 & ja \\
SF7 & ramp\_\allowbreak{}56\_\allowbreak{}down & - & - & 0 & ja \\
SF7 & peak & - & - & 1 & nein \\
SF8 & scaled & - & - & 0 & ja \\
SF8 & scaled\_\allowbreak{}down & - & - & 0 & ja \\
SF8 & ramp\_\allowbreak{}28 & - & - & 0 & ja \\
SF8 & ramp\_\allowbreak{}28\_\allowbreak{}down & 57 & weekly + long\_\allowbreak{}term & 0 & ja \\
SF8 & ramp\_\allowbreak{}56 & - & - & 0 & ja \\
SF8 & ramp\_\allowbreak{}56\_\allowbreak{}down & 57 & weekly + long\_\allowbreak{}term & 0 & ja \\
SF8 & peak & 1 & weekday & 0 & ja \\
\end{tablang}

\subsubsection*{Niveauänderungen: erster Tag, an dem die Skalen anhalten}
\begin{tablang}{W{0.89} W{1.04} V{0.74} V{1.04} V{1.19} V{1.11}}{Stichtag & Szenario & weekly hält an ab Tag & long\_\allowbreak{}term hält an ab Tag & beide halten an ab Tag (Erkennung) & weekly davor abgerissen [mal]}
SF6 & scaled & 7 & 59 & - & 4 \\
SF6 & scaled\_\allowbreak{}down & 7 & 22 & 31 & 1 \\
SF6 & ramp\_\allowbreak{}28 & 7 & 50 & 50 & 2 \\
SF6 & ramp\_\allowbreak{}28\_\allowbreak{}down & 7 & 28 & 31 & 1 \\
SF6 & ramp\_\allowbreak{}56 & 7 & 57 & 57 & 2 \\
SF6 & ramp\_\allowbreak{}56\_\allowbreak{}down & 7 & 30 & - & 2 \\
SF7 & scaled & 6 & 15 & 15 & 0 \\
SF7 & scaled\_\allowbreak{}down & 35 & 38 & 38 & 0 \\
SF7 & ramp\_\allowbreak{}28 & 7 & 17 & 17 & 0 \\
SF7 & ramp\_\allowbreak{}28\_\allowbreak{}down & 8 & 41 & 41 & 2 \\
SF7 & ramp\_\allowbreak{}56 & 7 & 17 & 29 & 1 \\
SF7 & ramp\_\allowbreak{}56\_\allowbreak{}down & 8 & 18 & - & 3 \\
SF8 & scaled & 1 & 28 & - & 3 \\
SF8 & scaled\_\allowbreak{}down & 1 & 8 & - & 1 \\
SF8 & ramp\_\allowbreak{}28 & 1 & 13 & - & 2 \\
SF8 & ramp\_\allowbreak{}28\_\allowbreak{}down & 1 & 9 & 57 & 1 \\
SF8 & ramp\_\allowbreak{}56 & 1 & 10 & - & 1 \\
SF8 & ramp\_\allowbreak{}56\_\allowbreak{}down & 1 & 9 & 57 & 1 \\
\end{tablang}
